\subsection{Dimensión de la imagen de un morfismo.}
\label{IV.2.s5.4-fr}

\begin{proposition}[5.4.1]
\label{IV.2.s5.4.1-fr}
Sean $X$, $Y$ dos preesquemas localmente noetherianos y $f:X\to Y$ un
morfismo.
\begin{enumerate}
\item[\textup{(i)}] Si $f$ es cuasifinito, se tiene
$\dim(X)\leq\dim(\overline{f(X)})\leq\dim(Y)$.
\item[\textup{(ii)}] Si $f$ es sobreyectivo y abierto \textup{(resp.} sobreyectivo y cerrado\textup{)}, se
tiene $\dim(X)\geq\dim(Y)$.
\end{enumerate}
\end{proposition}

\begin{proof}
\begin{enumerate}
\item[(i)] Se puede reemplazar $f$ por $f_{\mathrm{red}}$ (\textbf{II}, 6.2.4), y por
tanto suponer que $X$ e $Y$ son reducidos; si $Z$ es el subpreesquema cerrado
reducido de $Y$ que tiene por espacio subyacente $\overline{f(X)}$, se tiene
entonces $f=j\circ g$, donde $j:Z\to Y$ es la inyección canónica y
$g:X\to Z$ es un morfismo cuasifinito (\textbf I, 5.2.2 y \textbf{II}, 6.2.4). Es posible, pues,
limitarse al caso en que $f(X)$ es denso en $Y$. Para todo $x\in X$,
$\mathcal O_x$ es entonces un $\mathcal O_{f(x)}$-módulo cuasifinito
(\textbf{II}, 6.2.2), y por consiguiente $\mathfrak m_{f(x)}\mathcal O_x$ es un ideal
de definición de $\mathcal O_x$ $(\mathbf{0}_{\mathrm I},7.4.4)$; pero se sabe
(\textbf 0, 16.3.10) que, si $A\to B$ es un homomorfismo local de anillos locales
noetherianos con la propiedad de que, si $\mathfrak m$ es el ideal maximal de
$A$, $\mathfrak mB$ es un ideal de definición de $B$, entonces se tiene
$\dim(B)\leq\dim(A)$; esto completa la demostración de (i) en virtud de
(5.1.4).

\item[(ii)] La definición de la dimensión (\textbf 0, 14.1.2) muestra que basta probar
que, para toda sucesión $(y_i)_{0\leq i\leq n}$ de elementos distintos de $Y$
tal que $y_i\in\overline{\{y_{i+1}\}}$ para $0\leq i\leq n-1$, existe una
sucesión $(x_i)_{0\leq i\leq n}$ de puntos de $X$ tal que
$x_i\in\overline{\{x_{i+1}\}}$ para $0\leq i\leq n-1$ y $f(x_i)=y_i$ para
todo $i$. Supongamos primero que $f$ es sobreyectivo y \emph{abierto} y demostremos
la existencia de los $x_i$ por inducción sobre $i$; la existencia de
$x_0\in X$ tal que $f(x_0)=y_0$ se sigue de que $f$ es sobreyectivo. Si se han
determinado los $x_i$ para $i\leq m$ de modo que se satisfaga $f(x_i)=y_i$
para $i\leq m$ y $x_i\in\overline{\{x_{i+1}\}}$ para $i<m$, se observa que,
puesto que $f$ es abierto y $y_{m+1}$ es una generización de $y_m$, existe
$x_{m+1}\in f^{-1}(y_{m+1})$ que es una generización de $x_m$ (1.10.3), y
puede proseguirse la inducción.

Supongamos ahora que $f$ es sobreyectivo y \emph{cerrado} y demostremos la existencia
de los $x_i$ por inducción \emph{descendente} sobre $i$; la existencia de $x_n$ tal
que $f(x_n)=y_n$ se sigue nuevamente de que $f$ es sobreyectivo. Si se han
determinado los $x_i$ de modo que satisfagan las condiciones requeridas para
$i>m$, se observa que $f(\overline{\{x_{m+1}\}})$ es la clausura de
$\{f(x_{m+1})\}=\{y_{m+1}\}$, puesto que $f$
\oldpage[IV-2]{94}
es cerrado (Bourbaki,
\emph{Top. gén.}, chap. I\textsuperscript{er}, 3\textsuperscript{e} éd.,
\S~5, n\textsuperscript{o}~4, prop.~9); existe, pues,
$x_m\in\overline{\{x_{m+1}\}}$ tal que $f(x_m)=y_m$, y puede proseguirse
también la inducción descendente.
\end{enumerate}
\end{proof}

